%%%%%%%%%%%%%%%%%%%%%%%%%%%%%%%%%%%%%%%%%%%%%%%%%%%%%
%%%%%	 		Course Outline- Template		%%%%%
%%%%%		  Prepaired by Dr. R.M. Silva		%%%%%
%%%%%			Department of Statistics		%%%%%
%%%%%		  Faculty of Applied Sciences		%%%%%
%%%%%		University of Sri Jayewardenepura   %%%%%
%%%%%				Dec. 09th, 2017				%%%%%
%%%%%%%%%%%%%%%%%%%%%%%%%%%%%%%%%%%%%%%%%%%%%%%%%%%%%


% Please place the SJPLogo in the same folder that you are going to run this Latex file

\documentclass[a4paper,12pt]{article}
\usepackage{graphicx}
\usepackage{enumitem}
\usepackage{setspace}
\usepackage{multicol} 
\usepackage{color}
\usepackage[
top    = 2.00cm,
bottom = 2.00cm,
left   = 2.00cm,
right  = 2.00cm]{geometry}
\usepackage{hyperref}
\hypersetup{pdfstartview={XYZ null null 1.00}}

\begin{document}

\begin{figure}[ht]
	\begin{center}
		\includegraphics[angle=0,scale=0.05]{SJPLogo.jpg}
	\end{center}
\end{figure}

\vspace{-1cm}

\begin{center}
	University of Sri Jayewardenepura\\
	Faculty of Applied Scineces \\
	Department of Statistics
\end{center}


\noindent\rule{17cm}{0.4pt} 	%This will create a horizontal line

\begin{multicols}{3}
	\noindent\textbf{Batch: 2016/2017}
	
	\columnbreak
	\noindent\textbf{Year: 2020 }
	
	\columnbreak
	\noindent\textbf{Semester: First Semester}
	
\end{multicols}

\noindent\rule{17cm}{0.4pt}	% horizontal line

\vspace{0.5cm}
\noindent\textbf{Course Unit: STA 326 2.0 Programming and Data Analysis with R}\\

\noindent\textbf{Type of the course unit: Core for special degree students/ Optional for others}\\

\noindent\textbf{Pre-Requisites:}\\
\noindent\textit{{ STA 114 2.0 Probability and Distribution Theory I, STA 123 2.0 Probability and Distribution Theory II, STA 124 1.5 Data Analysis I, STA 213 2.0 Statistical Inference, STA 226 1.5 Data Analysis II}}\\

\noindent\textbf{Workload:}\\
\noindent\textit{{ Minimum total expected workload to achieve the learning outcomes for this unit is 100 hours per semester typically comprising a mixture of scheduled learning activities, independent study and 26 hours of lectures. Independent study may include associated readings, assessment and preparation for scheduled activities.}}\\

\noindent\textbf{Course Objective(s):}
\begin{itemize}
	\setlength\itemsep{0.1mm}
	\item To introduce how to program efficiently in R.
	\item To provide an in-depth, more advanced coverage of data processing, visualisation and analysis in the R programming environment.
\end{itemize}
\noindent\textbf{Course Contents:}
\begin{enumerate}[label*=\arabic*.]
	\setlength\itemsep{0.1mm}
	\item R basics: Objects in R, Data types, Operations, Installing packages, Control structures, Piping
	\item Writing functions in R
	\item Data analysis with the tidyverse
	\begin{enumerate}[label*=\arabic*.]
	\item Data import and export
	\item Data wrangling: Tidy data principles, Reshaping data into tidy form, Data transformation
	\item Data visualization: The grammar of graphics
	\item Statistical	modelling	and	inference
	\item Communication: Dynamic reproducible reporting
	\end{enumerate}
\end{enumerate}

\newpage

\noindent\textbf{Learning Outcomes:}
At the end of this course, the student should be able to:
\begin{itemize}
	\setlength\itemsep{0.1mm}
\item navigate the R integrated development environment (IDE) R Studio
\item execute basic arithmetic operations in R
\item define data classes, object attriutes, data structures in R
\item describe differences in different data structures
\item write user-defined functions  to solve a given problem in R 
\item solve fundamental error problems and bugs in R programs
\item describe the principles of tidyverse programming
\item use Tidyverse packages in data science workflow 
\item import external data into R for data processing and statistical analysis.
\item organise complex, messy, data into the most convenient form for analysis or reporting.
\item create data graphics using the ggplot2 package
\item select effective visualisation and modelling approaches to understand relationships between variables, and make decisions with data
\item interpret the results of analysis and communicate these to a broad audience.
\item Evaluate several functions that perform the same task and explain which one is better and why?
\item Revise functions/programs to improve the readability and computational efficiency

\end{itemize}

\noindent\textbf{Method of Assessment:}
\begin{itemize}
	\setlength\itemsep{0.1mm}
	\item Continuous assessment: 20\%
	\item Final examination: 80\%
\end{itemize}

\noindent\textbf{Recommended Readings:}
\begin{itemize}
	\setlength\itemsep{0.1mm}
	\item Course website: Everything you want to know about the course, and everything you will need for the course (links to weekly reading, tutorials and lecture materials) will be posted at \textcolor{blue}{hellor.netlify.com}\\
	Author: Thiyanga Talagala
	\item Title: R for Data Science \\
		  Author(s): Hadley Wickham and Garrett Grolemund \\
		  Publisher: O'REILLY \\
		  This book is available online for free. Visit \textcolor{blue}{https://r4ds.had.co.nz/}
		\item Title: Advanced R \\
		  Author(s): Hadley Wickham \\
		  Publisher: Chapman \& Hall/CRC \\
		  This book is available online for free. Visit \textcolor{blue}{https://adv-r.hadley.nz/}
\end{itemize}

\noindent\textbf{Lecturer in Charge:}

\begin{itemize}
	\setlength\itemsep{0.1mm}
	\item[] Name: Dr Thiyanga  Talagala
	\item[] Email: ttalagala@sjp.ac.lk
\end{itemize}


\noindent\textbf{\textit{Note: The course can be dropped on or before 21 February 2020.}}

\end{document}
